\chapter{Introducción}

\section{Antecedentes}
En los últimos años, los Sistemas de Gestión de Aprendizaje (LMS) han trascendido su función inicial como simples repositorios de contenido para consolidarse como el núcleo operativo de la educación digital contemporánea~\cite{Supriyanto2025,YeseniaZavaleta-Sanchez2025}. Estas plataformas ofrecen la capacidad técnica de registrar meticulosamente cada interacción, clic y trayectoria de navegación de los estudiantes, generando volúmenes masivos de datos asíncronos. La trascendencia de este fenómeno radica en que, al procesarse de manera estructurada, esta información permite modelar perfiles estudiantiles basados directamente en su comportamiento real dentro del entorno virtual~\cite{YeseniaZavaleta-Sanchez2025}. Esto ha motivado a la comunidad investigadora e instituciones a desarrollar métodos para extraer valor de estos registros continuos con el fin de personalizar la enseñanza, identificar riesgos de bajo rendimiento y evaluar la efectividad del diseño instruccional, consolidando a los LMS como un recurso analítico primordial en el entorno educativo~\cite{Mercado2024,Setiadi2025}.

Para estructurar y analizar metodológicamente esta información, la literatura ha adoptado principalmente dos grandes paradigmas: la Minería de Datos Educativos (EDM) y la Minería de Procesos Educativos (EPM)~\cite{Semler2025}. Por un lado, diversos estudios han aplicado técnicas predictivas y de aprendizaje automático propias de la EDM para evaluar y anticipar el rendimiento académico a partir del historial de interacción del usuario~\cite{Mercado2024,Setiadi2025}. Por otro lado, la EPM ha demostrado una notable eficacia al extraer conocimiento de los registros de eventos (\textit{event logs}) para descubrir, auditar y mapear visualmente las secuencias y rutas de aprendizaje empíricas en entornos asíncronos~\cite{AlQaheri2022,Supriyanto2025}. En conjunto, estas investigaciones validan la factibilidad técnica de caracterizar el comportamiento en línea y aportan evidencia crucial para fundamentar mejoras didácticas~\cite{Semler2025}.

Sin embargo, a pesar de estos sólidos avances, persiste una brecha tecnológica crítica que restringe su adopción como una herramienta confiable para la gestión basada en evidencia. La revisión sistemática de la literatura advierte que la mayoría de los esfuerzos actuales operan de forma fragmentada y bajo un enfoque de análisis estático o por lotes (\textit{batch}); es decir, construyen perfiles retrospectivos que no se adaptan ni actualizan conforme evoluciona el comportamiento del alumno durante su ciclo formativo~\cite{Semler2025,YeseniaZavaleta-Sanchez2025}. Además, una proporción significativa de estos modelos predictivos opera como sistemas opacos o de ``caja negra'', arrojando métricas sin proporcionar una justificación estructural que el profesor pueda analizar o interpretar pedagógicamente~\cite{BarredoArrieta2020,Lee2025}. Si bien existen marcos teóricos recientes que proponen la integración del procesamiento de flujos continuos de datos con minería de procesos~\cite{YeseniaZavaleta-Sanchez2025}, aún no se ha consolidado una plataforma de software unificada que automatice este ciclo analítico de principio a fin, capaz de generar perfiles estudiantiles dinámicos, fundamentados en su interacción constante, y que aseguren una interpretabilidad algorítmica indispensable para orientar el diseño de intervenciones pedagógicas oportunas.

\section{Planteamiento del problema}
La literatura reciente en Minería de Datos Educativos (EDM) y Minería de Procesos Educativos (EPM) ha demostrado avances significativos en el análisis del comportamiento estudiantil~\cite{Semler2025,Supriyanto2025}. Existen investigaciones consolidadas que predicen el rendimiento del alumno o auditan su historial de navegación dentro de un LMS de forma retrospectiva~\cite{Mercado2024,Setiadi2025}. Sin embargo, la limitación inherente a estos enfoques radica en la prevalencia del análisis tradicional basado en datos estáticos por lotes; se carece de la integración de técnicas capaces de operar sobre flujos continuos de eventos (\textit{data streams}), donde el perfil computacional del estudiante se ajuste incremental y automáticamente conforme evolucionan sus interacciones en el LMS a lo largo del curso~\cite{Batool2021,YeseniaZavaleta-Sanchez2025}.

Por otro lado, para que el sector docente adopte estos sistemas analíticos, es indispensable garantizar su transparencia y usabilidad operativa~\cite{BarredoArrieta2020}. Aunque existen marcos teóricos sólidos orientados a la creación de modelos dinámicos~\cite{YeseniaZavaleta-Sanchez2025}, o aproximaciones que recurren a la Inteligencia Artificial Explicable (XAI) para interpretar predicciones estáticas~\cite{Desai2025,Lee2025}, aún no se ha construido una plataforma de software unificada que automatice este proceso analítico de principio a fin. Específicamente, el desafío pendiente en la literatura científica consiste en generar perfiles estudiantiles computacionalmente dinámicos que el educador pueda interpretar con facilidad y sin ambigüedad estructural, aun cuando dichos perfiles cambien y se actualicen continuamente a partir de eventos capturados en tiempo real.

\section{Preguntas de investigación}
Para abordar la problemática planteada, esta investigación se guía por los siguientes cuestionamientos:

\begin{itemize}
    \item ¿En qué medida la automatización concurrente de la minería de flujos de datos y la minería de procesos logra mantener tiempos de procesamiento de baja latencia ante el volumen masivo de eventos asíncronos generados en un LMS?
    \item ¿Cuál es el grado de interpretabilidad y utilidad práctica que ofrecen los perfiles de usuario estructurados a partir de rutas secuenciales continuas, al ser expuestos a la evaluación o criterios de toma de decisiones docentes?
    \item ¿Cuáles son las diferencias operativas y de fidelidad representativa al contrastar el modelado del comportamiento estudiantil basado en un paradigma de análisis estático por lotes (\textit{batch}) frente a la minería de flujos de datos continuos (\textit{data streams})?
\end{itemize}

\section{Objetivos}
La definición de objetivos de investigación fundamenta el alcance metodológico y la dirección técnica de este proyecto. Considerando la complejidad de procesar y extraer valor de los flujos continuos de información generados en los Sistemas de Gestión de Aprendizaje (LMS), resulta indispensable trazar metas precisas que guíen la transición de los fundamentos teóricos hacia una plataforma computacional aplicable. A continuación, se expone el objetivo general, el cual rige el propósito central de este trabajo, y los objetivos específicos que delinean las etapas arquitectónicas, analíticas y de validación necesarias para su implementación.

\subsection{Objetivo General}
Automatizar el perfilado dinámico de estudiantes a partir de sus interacciones en un LMS, mediante la puesta en operación de una plataforma de software que implemente un método basado en minería de flujos de datos y de procesos, para obtener perfiles educativos interpretables, consistentes y reproducibles.

\subsection{Objetivos Específicos}
\begin{itemize}
    \item Diseñar la arquitectura de datos para la extracción, limpieza y estructuración continua de los registros de interacción provenientes del LMS.
    \item Desarrollar el módulo analítico que integre y ejecute coordinadamente los algoritmos de minería de flujos de datos y minería de procesos para generar los modelos dinámicos de usuario.
    \item Evaluar el desempeño de la plataforma computacional midiendo su eficiencia en el tiempo de procesamiento, la consistencia algorítmica, la reproducibilidad del método y la interpretabilidad de los perfiles resultantes.
\end{itemize}

\section{Justificación}
La analítica del aprendizaje en los Sistemas de Gestión de Aprendizaje (LMS) ha avanzado primordialmente en la predicción retrospectiva y el análisis de datos estáticos~\cite{Mercado2024,Setiadi2025}. Sin embargo, existe un vacío tecnológico sustancial en la integración de técnicas que procesen flujos continuos de eventos para ajustar automáticamente los perfiles computacionales conforme evolucionan las interacciones del estudiante~\cite{YeseniaZavaleta-Sanchez2025}. Atender este problema resulta fundamental, ya que la literatura advierte que los esfuerzos analíticos actuales suelen generar predicciones estáticas que operan como sistemas opacos, es decir, modelos de ``caja negra'' cuyo funcionamiento interno carece de la capacidad de ser explicado o presentado en términos comprensibles para el ser humano~\cite{BarredoArrieta2020,Lee2025}. En consecuencia, para que estas herramientas logren una verdadera adopción en el sector docente, es indispensable garantizar su usabilidad e interpretabilidad computacional frente a flujos de datos dinámicos~\cite{Desai2025,Lee2025}.

Ante esta necesidad, la presente investigación aportará evidencia y la sistematización técnica para operar una plataforma computacional unificada que automatice el método de modelado dinámico de extremo a extremo. Este desarrollo contribuirá a superar la fragmentación detectada en la literatura de minería de procesos educativos~\cite{Semler2025} y a consolidar la convergencia metodológica entre la minería de flujos de datos y la minería de procesos~\cite{YeseniaZavaleta-Sanchez2025}, garantizando un tiempo de procesamiento eficiente y la reproducibilidad algorítmica. Los principales beneficiarios serán los educadores, a quienes el sistema ofrecerá criterios para la toma de decisiones pedagógicas basadas en evidencia confiable. De este modo, la automatización analítica permitirá identificar patrones tempranos y rutas de aprendizaje reales mediante perfiles dinámicos, estructurados y directamente comprensibles para el usuario final.

La ejecución de este proyecto es técnica y metodológicamente viable gracias a la existencia de fundamentos teóricos sólidos que fungen como núcleo estructurado para la automatización, tales como los marcos recientes que integran la minería de flujos y procesos en contextos educativos~\cite{YeseniaZavaleta-Sanchez2025}. Lejos de partir de un enfoque especulativo, la propuesta materializará en software una estructura teórica formal. Adicionalmente, la viabilidad algorítmica y de infraestructura está respaldada por trabajos afines que demuestran la eficacia de las arquitecturas orientadas a eventos en tiempo real~\cite{Kundavaram2025} y el agrupamiento evolutivo sobre flujos continuos de datos~\cite{Batool2021}. La articulación de estas bases comprobadas brindará el rigor necesario para garantizar que la transición del paradigma estático al dinámico sea consistente y validable empíricamente.

\section{Alcances y limitaciones}

En esta sección se establecen las fronteras operativas y conceptuales de la investigación. A continuación, se delimita el alcance del desarrollo tecnológico de la plataforma y se explicitan las restricciones inherentes al procesamiento de datos asíncronos generados dentro de los Sistemas de Gestión de Aprendizaje (LMS).

\subsection{Alcances}

El alcance principal de este proyecto comprende la sistematización técnica y automatización metodológica del modelado dinámico de usuarios. Específicamente, la investigación abarcará las siguientes metas operativas:

\begin{itemize}
    \item \textbf{Desarrollo de la infraestructura de datos:} Se diseñará e implementará una arquitectura orientada al procesamiento de flujos continuos (arquitectura tipo Kappa) para automatizar la extracción, limpieza y estructuración de los registros de interacción de los estudiantes desde el LMS~\cite{Kundavaram2025}.
    \item \textbf{Integración del módulo analítico:} Se construirá un sistema unificado de software que ejecute, de forma coordinada, algoritmos de agrupamiento de minería de flujos de datos~\cite{Batool2021} junto con técnicas de minería de procesos~\cite{AlQaheri2022,Supriyanto2025} para generar rutas de aprendizaje y perfiles estudiantiles que evolucionen en el tiempo~\cite{YeseniaZavaleta-Sanchez2025}.
    \item \textbf{Validación computacional:} El desempeño de la plataforma será evaluado utilizando métricas técnicas de eficiencia, midiendo la latencia en el tiempo de procesamiento de eventos y garantizando la reproducibilidad algorítmica del método.
    \item \textbf{Evaluación de transparencia:} Se validará el grado de interpretabilidad de los perfiles resultantes, asegurando que el modelo exponga justificaciones estructurales comprensibles orientadas a la toma de decisiones docentes~\cite{BarredoArrieta2020,Lee2025}.
\end{itemize}

\subsection{Limitaciones}

Aunque la plataforma automatizará el análisis de datos en tiempo real, el desarrollo y ejecución del proyecto están sujetos a restricciones técnicas y de frontera que deben considerarse:

\begin{itemize}
    \item \textbf{Dependencia de los registros del LMS (\textit{Logs}):} El sistema de modelado funciona estrictamente a partir de los datos que la plataforma educativa logra registrar y almacenar~\cite{Setiadi2025}. Si el LMS no captura ciertas interacciones o secuencias de navegación de los estudiantes debido a limitaciones de instrumentación técnica, el algoritmo no tendrá manera de procesarlas ni de incluirlas en el perfil dinámico.
    \item \textbf{Frontera técnica vs. pedagógica:} El proyecto está enfocado en la ciencia de datos y la ingeniería de software para construir y validar la plataforma tecnológica. El sistema entregará los perfiles interpretables a los profesores como una herramienta de apoyo analítico, pero la investigación no evaluará experimentalmente el impacto directo de las intervenciones didácticas que los docentes implementen en el rendimiento final de los alumnos.
    \item \textbf{Prioridad en la interpretabilidad sobre predicciones puras:} Existen modelos predictivos de caja negra que alcanzan altos coeficientes de precisión matemática pero carecen de transparencia operativa~\cite{BarredoArrieta2020,Lee2025}. Esta investigación se limita deliberadamente a la generación de modelos descriptivos, transparentes y estructurados en rutas de procesos; no se desarrollarán arquitecturas de aprendizaje profundo complejas si estas comprometen la interpretabilidad directa requerida por el educador.
\end{itemize}